\section*{The Propers: Themes and Movement}
\label{triptych:brief-synthesis:start}

\begingroup\footnotesize

\begin{massbox}{Governing thesis: grace opens a people to receive the risen Christ, confess him truthfully, and return every received good in worship and service}
The Mass moves from God's holy dwelling and strength through apostolic
transmission and the opening of ears and tongue to accepted service,
firstfruits, and remedy. Its repeated contrast is not between spirit and body
or hearing and action, but between a gift closed upon itself and grace made
fruitful in confession, praise, offering, and communion.
\end{massbox}

\begin{enumerate}[leftmargin=1.5em,itemsep=0.1em]
\item \textbf{A people is gathered:} God establishes concord, strength, and a place of holiness.
\item \textbf{The gospel is received:} apostolic testimony gives the content to which hearing is opened.
\item \textbf{Reception becomes confession:} Christ heals without spectacle and makes speech right.
\item \textbf{Gift becomes service:} praise, firstfruits, sacrifice, and sacramental remedy return grace toward God and neighbor.
\end{enumerate}

\noindent\emph{How to read the movement.} The stages overlap. Psalmic prayer
already speaks and offers; Paul's witness already begins with reception; the
Gospel binds hearing and speech in one healing; the altar's service asks to
receive help. The architecture below is source-grounded synthesis, not a claim
that a known historical compiler intended every relation.

\subsection*{The house God opens}

The Introit begins with divine agency: God is in his holy place, makes people
dwell together, and gives strength to his people. The complete Psalm~67
widens this from a private interior to God's public care for orphan, widow,
prisoner, Israel, and Zion. The Missal's \latin{unanimes} stresses concord,
although the Clementine reads \latin{unius moris}; that difference must not
be disguised by an invented English translation. Here holiness is not
withdrawal from embodied need. It is the dwelling from which the vulnerable
are protected and a people is strengthened.

The Collect immediately prevents this gathered people from measuring grace
by entitlement. God's abundant goodness exceeds both the merits and the
requests of suppliants. Mercy answers what conscience fears and adds what
prayer does not presume. Blessed Ildefonso Schuster calls this prayer a
compact liturgical theology: forgiveness is joined to the gift beyond
ordinary request, which he reads toward perfect charity and union with God.
His reading is later spiritual reception, not a recovered account of the
prayer's origin. The prayer itself establishes the asymmetry: supplication is
real, conscience is real, and mercy is greater than either merit or desire.

This opening controls the bodily and social images that follow. ``House''
cannot be reduced to a modern program, nor can abundance be promised as
material prosperity. The text instead places concord, strength, conscience,
and hope under the giver whose action precedes and surpasses the petitioners.

\subsection*{A gospel received before it is spoken}

First Corinthians~15 names the content of Christian hearing. Paul does not
offer a private insight: he makes known the gospel he preached, which the
Corinthians received, in which they stand, and through which they are saved if
they hold it fast. He then says that he delivered what he also received:
Christ died for sins, was buried, rose on the third day according to the
Scriptures, and appeared to named and numerous witnesses. The sequence joins
tradition and testimony without making faith a free-standing feeling.

Ancient exegesis presses several distinct aspects. Irenaeus and Tertullian use
the death, burial, and Resurrection against denials of Christ's real flesh.
Chrysostom emphasizes transmission, the evidentiary force of burial and
appearances, and Paul's humility. Theodoret reads the witnesses as a publicly
testable accumulation. Aquinas's reportatio on the chapter distinguishes
received doctrine, Scriptural fulfillment, and the several senses of Paul's
``untimely birth.'' These are not interchangeable proofs, but each refuses a
gospel detached from Christ's bodily death and Resurrection or from the
apostolic act of handing on.

The lection's exact boundary is decisive. It ends within v.~10 at grace that
``has not been void.'' It does not proclaim Paul's following statement that he
labored more than the others, ``yet not I, but the grace of God with me.''
Augustine, Gregory, Aquinas, and later Catholic teaching make important use of
that full canonical sentence for grace and human cooperation. Their doctrine
may illuminate the chapter, but the omitted clause may not be presented as
though heard in this Mass. The appointed ending instead leaves the transformed
persecutor and grace's fruitfulness as its last note.

\newpage

\subsection*{Opened ears and right speech}

Mark's scene supplies neither a generic symbol nor an abstraction from the
man's body. Unnamed people bring a particular person with deafness and impaired
speech; Jesus takes him apart, touches ears and tongue, looks to heaven,
groans, and speaks a Semitic imperative conventionally identified as Aramaic,
whose sound the evangelist preserves: \latin{Ephphetha}. Hearing opens, the bond of the tongue is released, and the
man speaks rightly. The body's restoration is itself a good. A responsible
spiritual reading therefore never equates disability with unbelief or moral
fault.

The direct tradition is unusually concrete. Gregory the Great, in
\emph{Homilies on Ezekiel} I.10.20, interprets the fingers as gifts of the
Holy Spirit, the Lord's saliva as wisdom flowing from divine speech, the sigh
as an example for penitential longing, and right speech as teaching supported
by prior action. Under the 1960 Office this very paragraph was appointed at
Matins for the Eleventh Sunday, making it a received liturgical reading of the
Sunday's Gospel rather than merely a remote parallel. Bede develops much the
same sacramental and catechetical field: Christ separates the sufferer from
the crowd, opens the ears of the heart through the Spirit, gives the savor of
wisdom, and forms confession. Theophylact stresses the holiness of Christ's
assumed flesh and his compassion for human nature. None licenses contempt for
the body; all begin from Christ's bodily act.

The command to silence complicates any simple ``opened mouth means speak
more'' lesson. Augustine asks why Christ forbids publication when he knows it
will increase, and reads the paradox as a rebuke to authorized preachers who
are sluggish. A 2025 catechesis of Pope Leo~XIV places the command within
Mark's paschal pedagogy: confession cannot yet be adequate until the whole
journey to the Cross and Resurrection is known. Hearing therefore governs
speech. The restored tongue is not permission to say anything; right speech
answers a received reality and its proper time.

\subsection*{From healing to service and firstfruits}

The Gradual and Offertory keep the miracle inside prayer. A heart trusts and
is helped; flesh flourishes again; the petitioner begs God not to be silent
or depart. Later the healed singer extols the Lord who upheld and healed him.
The Alleluia turns help into festival praise. These chants do not dissolve
Mark's unique action into a repeated technique. They give the assembly words
for dependence and thanksgiving before and after the Gospel.

At the altar, \latin{servitus}, \latin{munus}, and \latin{subsidium} bind
action and reception. The Secret names what the Church offers as service, yet
asks God to make the gift acceptable and to make it help for fragility.
Schuster hears in \latin{servitus} the priestly \emph{leitourgia} of the
Church and in the Eucharist strength for weakness. The Communion then places
substance and firstfruits under the honor of God. Proverbs' barns and wine
presses belong to sapiential instruction in trust and generosity; they are
not a guarantee that every giver will become wealthy. Catholic reading is
also governed by Christ's rejection of anxious accumulation and by the
Church's preferential care for those in need.

The Postcommunion completes the movement without dividing the human person.
Sacramental reception is asked to bring help to mind and body so that both may
be saved and the recipients may glory in the fullness of heavenly remedy.
The prayer neither reduces sacramental fruit to present bodily cure nor treats
the body as disposable. Schuster relates contact with Christ's Body to present
strength and final resurrection. Thus the house, risen body, healed senses,
flourishing flesh, offered substance, and whole-person remedy converge on one
bounded claim: grace received becomes fruitful when it opens persons to true
confession, worship, service, and the hope of bodily salvation.

\begin{threestable}{Movement}{Source-grounded fruit}{Controlling limit}
\textbf{Gathered and strengthened} & Divine mercy establishes a concordant people and exceeds merit and request. & Neither concord nor abundance is a warrant for coercion, exclusion, or prosperity teaching.\\
\textbf{Receiving and handing on} & The Resurrection gospel precedes and governs Christian confession. & The lection stops within v.~10; later exegesis of the omitted clause remains canonical context.\\
\textbf{Hearing and speaking rightly} & Christ's concrete healing supports baptismal, catechetical, and moral reception. & Physical disability is not a figure for culpability, and secrecy qualifies indiscriminate speech.\\
\textbf{Offering and receiving help} & Service, firstfruits, Eucharistic remedy, mind, and body belong to one responsive movement. & Liturgical synthesis does not promise automatic wealth or temporal healing.\\
\end{threestable}

\endgroup

\label{triptych:brief-synthesis:end}
\clearpage
